% Work/MAT -- 1. rocnik -- hodina 38 -- finalni revize 2026-10-30.
% Priklady z uzivatelem dodanych obrazku upraveny a doplneny vlastnimi ulohami.
% Vysledky a podminky kontrolovany nezavisle; zapis vysledku nepredpoklada
% kladnost promennych tam, kde to neni stanoveno.
\documentclass[11pt,a4paper]{article}
% Spolecny styl pro ucitelske pripravy MAT (Work/MAT)
% Verze 1.1, 2026-10-23
\NeedsTeXFormat{LaTeX2e}
\ProvidesPackage{priprava}[2026/10/23 v1.1 Sablona ucitelskych priprav MAT]

% Zakladni typografie a cestina
\RequirePackage[T1]{fontenc}
\RequirePackage[utf8]{inputenc}
\RequirePackage[czech]{babel}
\RequirePackage{lmodern}
\RequirePackage{microtype}
\RequirePackage{amsmath,amssymb}
\RequirePackage{xcolor}
\RequirePackage{enumitem}
\RequirePackage[most]{tcolorbox}
\RequirePackage{tabularx,booktabs,array}
\RequirePackage{fancyhdr}
\RequirePackage{needspace}
\RequirePackage{geometry}
\RequirePackage[hidelinks]{hyperref}

\geometry{a4paper,top=17mm,bottom=18mm,left=23mm,right=23mm}
\definecolor{notegray}{gray}{0.39}
\setlength{\parindent}{0pt}
\setlength{\parskip}{2.5pt plus 1pt minus 0.5pt}
\setlist[itemize]{topsep=2pt,itemsep=1pt,parsep=0pt,leftmargin=1.45em}
\setlist[enumerate]{topsep=2pt,itemsep=2pt,parsep=0pt,leftmargin=1.65em}
\setcounter{secnumdepth}{0}
\renewcommand{\arraystretch}{1.12}

% Zahlavi: vlevo rocnik | globalni tema, vpravo cislo: nazev hodiny
% Zapati zustava zcela prazdne.
\pagestyle{fancy}
\fancyhf{}
\renewcommand{\headrulewidth}{0.35pt}
\renewcommand{\footrulewidth}{0pt}
\setlength{\headheight}{14pt}
\newcommand{\nastavPripravu}[4]{%
  \fancyhead[L]{\small #1\enspace\textbar\enspace #2}%
  \fancyhead[R]{\small #3: #4}%
}

% Minimalisticke cernobile boxy. Kazda definice zvlast.
\tcbset{commonlesson/.style={enhanced,breakable=false,colback=white,colframe=black,
  colbacktitle=white,coltitle=black,boxrule=0.45pt,
  arc=0.8mm,outer arc=0.8mm,left=3mm,right=3mm,top=1.6mm,bottom=1.6mm,
  toptitle=1mm,bottomtitle=0.7mm,fonttitle=\bfseries,
  attach boxed title to top left={xshift=2.5mm,yshift=-1.15mm},
  boxed title style={colback=white,colframe=white,boxrule=0pt,left=1mm,right=1mm}}}
\newtcolorbox{zadaniBox}[1][]{commonlesson,title={Zadání},#1}
\newtcolorbox{prikladBox}[1][]{commonlesson,title={Příklad},#1}
\newtcolorbox{definiceBox}[1][]{commonlesson,title={Definice},#1}

% Spolecny box pro zadani a jeho reseni, bez mezery mezi castmi.
% Pouziti:
%   \begin{prikladReseniBox}
%     Zadani ulohy.
%     \tcblower
%     \begin{align*} ... \end{align*}
%   \end{prikladReseniBox}
% Spodni cast ma jemny sedy podklad a delici linku; lze v ni sazet
% vice radku matematiky. Obe casti maji jen jeden spolecny ramecek.
\tcbset{solutionlesson/.style={
  commonlesson,
  enhanced jigsaw,
  breakable,
  colbacklower=black!3,
  segmentation style={solid,draw=black!25,line width=0.35pt},
  fontlower=\small\color{notegray},
  before lower={\noindent\textbf{Řešení:}\enspace},
  leftlower=3mm,rightlower=3mm
}}
\newtcolorbox{prikladReseniBox}[1][]{solutionlesson,title={Příklad},#1}
\newtcolorbox{zadaniReseniBox}[1][]{solutionlesson,title={Zadání},#1}

% Pomucky pro autora pripravy: sedive poznamky; \reseni pro starsi pripravy
\newcommand{\poznamka}[1]{\par\begingroup\small\color{notegray}\itshape
  \textnormal{\textbf{Poznámka pro učitele.}}\enspace #1\par\endgroup}
\newcommand{\reseni}[1]{\par\begingroup\small\color{notegray}
  \textbf{Očekávané řešení:}\enspace #1\par\endgroup}
\newcommand{\blok}[1]{\needspace{5\baselineskip}\subsection*{#1}}
\newcommand{\vyrok}[1]{\uv{#1}}

\nastavPripravu{1. ročník}{Mocniny a odmocniny}{38}{KISS II: velké opakování}
\begin{document}

\poznamka{\textbf{Záměr:} nejdříve nahlas pojmenuj vhodnou strategii,
teprve potom počítej. Povinná část na 45 minut: 5 min připomenutí KISS,
9 min Příklad 1 + vybrané položky Samostatné práce 1,
10 min Příklad 2 + Samostatná práce 2 a), b),
10 min Příklad 3 + Samostatná práce 3 a), c),
8 min Příklad 4 a Samostatná práce 4 a), b),
3 min závěr a zadání domácí práce. Zbytek je cíleně \textbf{zásobník náročných úloh},
nikoli očekávání, že všichni vyřeší celý dokument.}

\blok{Připomenutí strategie: KISS = Keep It Simple, Stupid}
Ve složeném výrazu není vítězem ten, kdo nejrychleji počítá obrovská čísla,
ale ten, kdo \textbf{najde nejjednodušší správný postup}.
\begin{definiceBox}[title={Strategie zjednodušování}]
\begin{enumerate}[label=\arabic*)]
\item \textbf{Zkontroluj podmínky} původního výrazu (jmenovatel, záporné mocniny, sudá odmocnina).
\item \textbf{Sjednoť zápis}: stejný základ, vhodná mocnina prvočísla nebo podobné odmocniny.
\item \textbf{Zmenšuj výraz}: nejprve upravuj \emph{uvnitř}, potom vnější mocniny a velké součiny.
\item \textbf{Každý řádek dělej jednu věc}. Znaménko mínus a závorky neignoruj.
\item \textbf{Výsledek zkontroluj}: znaménko, exponenty, velikost, původní podmínky.
\end{enumerate}
\end{definiceBox}
\poznamka{Zkratku KISS představ jako neformální programátorskou zásadu jednoduchého řešení;
v matematice musí každá úprava současně zachovat hodnotu na původním definičním oboru.}

\blok{1. Algebraické mocniny: podíl zlomků, znaménka a podmínky}
\begin{prikladReseniBox}[title={Příklad 1: nejprve dělení zlomků},breakable=false]
Zjednoduš a uveď podmínky:
\[
\text{a) }\frac{a^2b}{c^3}:\frac{ab^2}{c^2},\qquad
\text{b) }\frac{a^2}{c^3}:\left(\frac a{c^2}\right)^{-2}.
\]
\tcblower
Pro všechna uvedená $a,b,c\ne0$ (v b) pouze $a,c\ne0$) platí
\[
\text{a) }\frac{a^2b}{c^3}\cdot\frac{c^2}{ab^2}=\boxed{\frac a{bc}},
\qquad
\text{b) }\frac{a^2}{c^3}\cdot\left(\frac a{c^2}\right)^2
=\boxed{\frac{a^4}{c^7}}.
\]
Záporný exponent převrací \emph{celý} základ, ne pouze část zlomku.
\end{prikladReseniBox}

\begin{zadaniReseniBox}[title={Samostatná práce 1: výrazy z vyšší úrovně},breakable=false]
Uprav (kde se vyskytuje podíl, urči jeho původní podmínky).
\begin{enumerate}[label=\alph*)]
\item $\displaystyle\frac{a^{-3}b^2}{c^{-3}}:\left(\frac{b^3c}{a^2}\right)^2$.
\item $\displaystyle\left(\frac{9a^{-2}c^3}{8b^{-2}}\right)^{-3}$.
\item $\displaystyle\frac{(-2)^5(-2)^{-3}}{4^{-2}}$.
\item $\displaystyle\left(\frac{a^2b^3}{a^3b^2}\right)^3:\left(\frac ba\right)^2$.
\item $(x-y)^3(y-x)^4$.
\end{enumerate}
\tcblower
\begin{enumerate}[label=\alph*)]
\item $\boxed{ac/b^4}$; $a,b,c\ne0$.
\item $\boxed{512a^6/(729b^6c^9)}$; $a,b,c\ne0$.
\item $(-2)^{5-3}\cdot4^2=4\cdot16=\boxed{64}$.
\item $\boxed{b/a}$; $a,b\ne0$.
\item $\boxed{(x-y)^7}$ (protože $(y-x)^4=(x-y)^4$); bez omezení.
\end{enumerate}
\end{zadaniReseniBox}

\begin{zadaniReseniBox}[title={Rozšiřující příklad 1: rozpoznáš nejkratší strategii?},breakable=false]
\[
\text{a) }\frac{3(a^3b^2)^3}{(2a^2b)^2}
\cdot\left(\frac{2a^4b^3}{a^3b}\right)^3,
\qquad
\text{b) }\left(\frac{2a^5b^4}{3a^3b}\right)^3:
\frac{(2a^3b^2)^3}{9(2a^2b)^2}.
\]
\tcblower
Pro $a,b\ne0$:
\[
\text{a) }\frac{3a^9b^6}{4a^4b^2}\cdot(2ab^2)^3
=\boxed{6a^8b^{10}},\qquad
\text{b) }\frac{8a^6b^9}{27}\cdot\frac{9}{2a^5b^4}
=\boxed{\frac{4ab^5}{3}}.
\]
\end{zadaniReseniBox}

\blok{2. Číselné mocniny: rozklad na prvočísla místo velkých čísel}
\begin{prikladReseniBox}[title={Příklad 2: zjednodušuj na prvočíslech},breakable=false]
Vyjádři součin mocnin prvočísel a pak číslo:
\[
\frac{8^2\cdot6^3}{4^3\cdot12}.
\]
\tcblower
\[
\frac{2^6\cdot2^3\cdot3^3}{2^6\cdot2^2\cdot3}
=2^{9-8}3^{3-1}=2\cdot3^2=\boxed{18}.
\]
\end{prikladReseniBox}

\begin{zadaniReseniBox}[title={Samostatná práce 2: bez vyčíslování obrovských mocnin},breakable=false]
\begin{enumerate}[label=\alph*)]
\item $\displaystyle\frac{125^2\cdot2^{10}}{16^2\cdot50}\cdot\frac{25}{5^3}$.
\item $\displaystyle\frac{3^7\cdot9^{-2}}{27^{-1}}$.
\item $\displaystyle\left(\frac{2^2}{5}\right)^5\left(-\frac{5^2}{2^3}\right)^3$.
\item $\displaystyle\frac{(72\cdot25)^3\cdot2^4}{5\cdot15^4\cdot32^2}$.
\end{enumerate}
\tcblower
\begin{enumerate}[label=\alph*)]
\item $2^{10-8-1}\cdot5^{6-2+2-3}=2\cdot5^3=\boxed{250}$.
\item $3^{7-4+3}=\boxed{3^6}$.
\item $-2^{10-9}\cdot5^{6-5}=\boxed{-10}$.
\item $2^{13-10}3^{6-4}5^{6-5}=2^3\cdot3^2\cdot5=\boxed{360}$.
\end{enumerate}
\end{zadaniReseniBox}

\begin{zadaniReseniBox}[title={Rozšiřující příklad 2: dlouhý výraz, krátký výsledek},breakable=false]
Bez zlomkové čáry ve výsledku, pouze pomocí \emph{celočíselných mocnin prvočísel},
uprav
\[
\left(\frac{8\cdot10^2\cdot9^{-3}}{15^4\cdot12^{-3}}\right)^{-3}
:\left(\frac{24^{-1}\cdot6^3}{20^2\cdot5^{-4}}\right)^2.
\]
\tcblower
První závorka $=2^{11}3^{-7}5^{-2}$, druhá $=2^{-4}3^25^2$.
Po umocnění a dělení:
\[
\frac{2^{-33}3^{21}5^6}{2^{-8}3^45^4}
=\boxed{2^{-25}\cdot3^{17}\cdot5^2}.
\]
\end{zadaniReseniBox}

\blok{3. Odmocniny, usměrňování a kontrola znamének}
\begin{prikladReseniBox}[title={Příklad 3: nejprve vyjmi čtverce},breakable=false]
Zjednoduš:
\[
\text{a) }\frac{(\sqrt{72}-\sqrt{18})^2}{(\sqrt2)^2},\qquad
\text{b) }\frac1{3-2\sqrt2}.
\]
\tcblower
\[
\text{a) }\frac{(6\sqrt2-3\sqrt2)^2}{2}=\boxed9,
\qquad
\text{b) }\frac{3+2\sqrt2}{9-8}=\boxed{3+2\sqrt2}.
\]
\end{prikladReseniBox}

\begin{zadaniReseniBox}[title={Samostatná práce 3: špatná strategie přidá práci},breakable=false]
\begin{enumerate}[label=\alph*)]
\item $(\sqrt{75}-\sqrt{12})(\sqrt3+1)$.
\item $\displaystyle\frac{\sqrt{48}+\sqrt{27}}{\sqrt3}$.
\item $\displaystyle\frac1{\sqrt7+\sqrt6}+\frac1{\sqrt7-\sqrt6}$.
\item $\displaystyle\sqrt[3]{-432}:\sqrt[3]{2}$.
\item Pro $t\in\mathbb R$ uprav $\sqrt{(t-4)^2}$.
\end{enumerate}
\tcblower
\begin{enumerate}[label=\alph*)]
\item $3\sqrt3(\sqrt3+1)=\boxed{9+3\sqrt3}$.
\item $(4\sqrt3+3\sqrt3)/\sqrt3=\boxed7$.
\item $(\sqrt7-\sqrt6)+(\sqrt7+\sqrt6)=\boxed{2\sqrt7}$.
\item $\sqrt[3]{-216}=\boxed{-6}$.
\item $\boxed{|t-4|}$, protože sudá odmocnina není záporná.
\end{enumerate}
\end{zadaniReseniBox}

\blok{4. Kombinace s technickým zápisem a jednotkami}
\begin{prikladReseniBox}[title={Příklad 4: odhad řádu před výpočtem},breakable=false]
Závod vyrobí $2{,}4\cdot10^5$ součástek za den.
Kolik jich vyrobí za $25$ dnů? Výsledek uveď vědeckým zápisem.
\tcblower
$2{,}4\cdot25\cdot10^5=60\cdot10^5=\boxed{6\cdot10^6}$ součástek.
\end{prikladReseniBox}

\begin{zadaniReseniBox}[title={Samostatná práce 4: exponentská a technická rutina},breakable=false]
\begin{enumerate}[label=\alph*)]
\item $\displaystyle\frac{8{,}4\cdot10^{-5}}{2{,}1\cdot10^{-6}}$.
\item $\displaystyle(6\cdot10^7)(1{,}5\cdot10^{-4})$.
\item Snímač uloží šest souborů po $256\,\mathrm{KiB}$. Kolik je to $\mathrm{MiB}$?
\item Čtvercová deska má obsah $288\,\mathrm{cm}^2$. Urči délku strany přesně.
\end{enumerate}
\tcblower
\begin{enumerate}[label=\alph*)]
\item $\boxed{40}$.
\item $\boxed{9\cdot10^3}$.
\item $6\cdot256/1024=\boxed{1{,}5\,\mathrm{MiB}}$.
\item $a=\sqrt{288}=\boxed{12\sqrt2\,\mathrm{cm}}$.
\end{enumerate}
\end{zadaniReseniBox}

\blok{Domácí úkol: dvacet příkladů ve čtyřech úrovních}
\poznamka{Záměrně velká nabídka: zadání lze rozdělit na povinných 10 položek
a 10 dobrovolných výzev. Požaduj celý postup a podmínky tam, kde jsou
proměnné ve jmenovateli nebo záporném exponentu.}

\begin{zadaniReseniBox}[title={Domácí úkol 1: mocniny a znaménka (1--5)},breakable=false]
\begin{enumerate}[label=\arabic*)]
\item $\displaystyle\frac{2^{-4}\cdot8^3}{4^2}$.
\item $\displaystyle\frac{((-3)^2)^3}{3^4}$.
\item $(-2)^5-(-2)^4$.
\item $\displaystyle\left(\frac35\right)^{-2}$.
\item $27^{2/3}$.
\end{enumerate}
\tcblower
\begin{enumerate}[label=\arabic*)]
\item $2^{-4+9-4}=\boxed2$.
\item $3^{6-4}=\boxed9$.
\item $-32-16=\boxed{-48}$.
\item $\boxed{25/9}$.
\item $\boxed9$.
\end{enumerate}
\end{zadaniReseniBox}

\begin{zadaniReseniBox}[title={Domácí úkol 2: proměnné a podmínky (6--10)},breakable=false]
\begin{enumerate}[label=\arabic*),start=6]
\item $81^{-3/4}$.
\item $\displaystyle\frac{a^7b^{-2}}{a^{-3}b^4}$.
\item $\displaystyle\left(\frac{2a^{-1}b^2}{3c^{-2}}\right)^{-2}$.
\item $\displaystyle\left(\frac{a^3b^2}{ab^{-1}}\right)^2$.
\item $(x-y)^2(y-x)^3$.
\end{enumerate}
\tcblower
\begin{enumerate}[label=\arabic*),start=6]
\item $\boxed{1/27}$.
\item $\boxed{a^{10}/b^6}$; $a,b\ne0$.
\item $\boxed{9a^2/(4b^4c^4)}$; $a,b,c\ne0$.
\item $\boxed{a^4b^6}$; $a,b\ne0$.
\item $\boxed{-(x-y)^5}$; bez omezení.
\end{enumerate}
\end{zadaniReseniBox}

\begin{zadaniReseniBox}[title={Domácí úkol 3: odmocniny a zlomky (11--15)},breakable=false]
\begin{enumerate}[label=\arabic*),start=11]
\item $\sqrt{200}-2\sqrt{18}+\sqrt8$.
\item $3\sqrt{75}-\sqrt{12}$.
\item $\sqrt[3]{-250}+2\sqrt[3]{16}$.
\item Pro $t\in\mathbb R$ uprav $\sqrt{(t-3)^2}$.
\item $\displaystyle\frac1{\sqrt5-2}$.
\end{enumerate}
\tcblower
\begin{enumerate}[label=\arabic*),start=11]
\item $\boxed{6\sqrt2}$.
\item $\boxed{13\sqrt3}$.
\item $\boxed{-\sqrt[3]2}$.
\item $\boxed{|t-3|}$.
\item $\boxed{\sqrt5+2}$.
\end{enumerate}
\end{zadaniReseniBox}

\begin{zadaniReseniBox}[title={Domácí úkol 4: kombinace a technika (16--20)},breakable=false]
\begin{enumerate}[label=\arabic*),start=16]
\item $\displaystyle\frac2{\sqrt7+\sqrt3}$.
\item $\displaystyle\frac{\sqrt{18}-\sqrt8}{\sqrt2}$.
\item $\displaystyle\frac{9\cdot10^5}{3\cdot10^{-2}}$.
\item $0{,}0000048\,\mathrm m$ převeď na $\mathrm{\mu m}$.
\item $\displaystyle\frac{125^2\cdot2^{10}}{16^2\cdot50}\cdot\frac{25}{5^3}$.
\end{enumerate}
\tcblower
\begin{enumerate}[label=\arabic*),start=16]
\item $\boxed{(\sqrt7-\sqrt3)/2}$.
\item $\boxed1$.
\item $\boxed{3\cdot10^7}$.
\item $\boxed{4{,}8\,\mathrm{\mu m}}$.
\item $\boxed{250}$.
\end{enumerate}
\end{zadaniReseniBox}
\end{document}
